% Copyright 2026 The terCAD team. All rights reserved.
% Use of this content is governed by a CC BY-NC-ND 4.0 license that can be found in the LICENSE file.

\subsection{Laravel}
\markboth{Bootcamping}{Laravel} \label{laravel}

Laravel is an open-source PHP framework created by Taylor Otwell in 2011 for building web
applications with expressive and elegant syntax. Following the Model-View-Controller (MVC)
pattern, it ships with everything required for a typical application: routing, templating,
database access, authentication, caching, queues, scheduling, and testing. The framework favors
convention over configuration, so most features work without any setup, while every part can be
replaced or extended when needed.

\paragraph{Request Lifecycle} Every request enters through \q{public/index.php}, which boots the
application from \q{bootstrap/app.php}. The HTTP kernel passes the request through a pipeline of
middleware. The router matches the request to a controller (or closure), and the returned value
is converted into a response. This cycle repeats for each request on a server.

\paragraph{Service Container} is the core of the framework that manages class dependencies and
performs dependency injection. Classes declared as arguments of constructors and methods are
resolved automatically. Service providers register bindings in the container and bootstrap
components (routes, events, and plugins) when the application starts up.

\paragraph{Artisan} is a command-line interface that provides helpers for development:
generating classes (\q{php artisan make:controller}), running migrations
(\q{php artisan migrate}), clearing caches (\q{php artisan optimize:clear}),
starting a REPL (\q{php artisan tinker}), and executing tests (\q{php artisan test}).

\paragraph{Ecosystem} Laravel is surrounded by a rich set of first-party packages: \q{Sanctum}
and \q{Passport} (API authentication), \q{Horizon} (queue monitoring), \q{Reverb} (WebSockets),
\q{Pint} (code style), \q{Sail} (Docker environment), and \q{Pail} (log tailing). Front-end
layers are covered by \q{Livewire} (reactive components without writing JavaScript) and
\q{Inertia} (single-page applications with Vue / React), both used by the official starter
kits.


\subsubsection{Structuring Application}

Laravel defines a conventional directory layout, where each component has its own place:

\dirtree{%
.1 /.
.2 app/\DTcomment{application code}.
.3 Console/Commands/\DTcomment{Artisan commands}.
.3 Http/Controllers/\DTcomment{request handlers}.
.3 Http/Middleware/\DTcomment{request filters}.
.3 Jobs/\DTcomment{queued tasks}.
.3 Models/\DTcomment{Eloquent models}.
.3 Providers/\DTcomment{service providers}.
.2 bootstrap/app.php\DTcomment{application configuration}.
.2 config/\DTcomment{configuration files}.
.2 database/migrations/\DTcomment{schema changes}.
.2 resources/views/\DTcomment{Blade templates}.
.2 routes/web.php\DTcomment{routes declaration}.
.2 storage/\DTcomment{logs, cache, files}.
.2 tests/\DTcomment{Unit and Feature tests}.
}

\noindent Environment-specific values (credentials, application id, debug mode) are kept in
\q{.env}-file and read via \q{env()}-helper inside \q{config/*}-files, while the code accesses
them through \q{config('app.name')}.


\subsubsection{Routing Requests}

Routes map URLs and HTTP methods to controllers. They can be grouped, named, and wrapped by
middleware:

\begin{lstlisting}
use App\Http\Controllers\ForumController;
use Illuminate\Support\Facades\Route;

Route::get('/cgi/forum_rooms.php',
    [ForumController::class, 'index']);

Route::post('/cgi/forum.php',
    [ForumController::class, 'room']);

// Parameters, names, and groups
Route::get('/hero/{id}',
    [HeroController::class, 'show'])
    ->name('hero');

Route::middleware('auth')->prefix('admin')->group(
  function () {
    Route::get('/stats', fn() => view('stats'));
  });
\end{lstlisting}

\noindent Middleware filters requests before (and after) they reach a controller. Since Laravel
11, they are registered in \q{bootstrap/app.php}:

\begin{lstlisting}
->withMiddleware(function (Middleware $middleware): void {
  $middleware->web([HotfixForPostAction::class, EnsureAgeAssurance::class]);
})
\end{lstlisting}

\begin{lstlisting}
class EnsureAgeAssurance
{
  private const EXEMPT = ['age', 'shell/*', 'up'];

  public function handle(Request $request, Closure $next): Response
  {
    if ($request->is(...self::EXEMPT)) {
      return $next($request);
    }
    if (!AgeAssurance::isDeclared()) {
      return redirect('/age');
    }
    if (!AgeAssurance::isAdult()) {
      return response()->view('restricted', [
        'isAdult' => false,
        'adultAge' => AgeAssurance::ADULT_AGE,
      ], 403);
    }
    return $next($request);
  }
}
\end{lstlisting}


\subsubsection{Handling Controllers}

Controllers group related request-handling logic. Input data is taken from the
\q{Request}-object (or \q{request()}-helper), and validated by rules before usage (our
\q{ForumController} extended by validation):

\begin{lstlisting}
final class ForumController extends Controller
{
  protected ForumParser $parser;

  public function __construct() {
    parent::__construct();
    $this->parser = new ForumParser();
  }

  public function topic(Request $request): View
  {
    $a = $request->validate([
      'id' => 'required|integer',
      'p' => 'nullable|integer|min:1', // page
    ]);
    $html = $this->get('cgi/f_show_thread.php', $a);
    $w = $this->get('cgi/w.JS', []);
    return view('forum_topic',
        $this->parser->parseTopic($html, $w));
  }
}
\end{lstlisting}

\noindent The returned value might be a view (\q{view('forum\_topic', \$arr)}),
a redirect (\q{redirect('/login')}), or JSON (\q{response()->json(\$data)}).

\subsubsection{Templating with Blade}

Blade is a templating engine compiled into plain PHP and cached. It supports layouts,
components, and directives:

\begin{lstlisting}[language=HTML]
{{-- resources/views/topics.blade.php (example) --}}
@extends('layouts.app')

@section('content')
  <h1>{{ $title }}</h1>  {{-- escaped output --}}
  {!! $html !!}           {{-- raw output --}}
  @forelse ($topics as $topic)
    <x-topic-item :topic="$topic" />
  @empty
    <p>No topics yet</p>
  @endforelse
  @if ($isGuest) @include('partials.login') @endif
@endsection
\end{lstlisting}

\noindent Anonymous components are placed into \q{resources/views/components} and called by
\q{<x-name>}-tag, while \q{@csrf}, \q{@auth}, \q{@env}, and \q{@vite} cover common needs
(security token, authorization, environment checks, and assets bundled by Vite).


\subsubsection{Persisting Data}

The Eloquent ORM represents each table as a model class. The schema is described by migrations,
which makes database changes tracked and repeatable:

\begin{lstlisting}
// database/migrations/2026_..._table.php
Schema::create('notifications', function (Blueprint $table) {
  $table->id();
  $table->timestamps();
  $table->longText('message');(*@ \stopnumber @*)
});

// database/migrations/2026_...._table.php
Schema::table('notifications', function (Blueprint $table) {
  $table->integer('channel')->default(0);(*@ \stopnumber @*)
});

// Usage of the Notification model
Notification::create([
    'message' => 'Hello',
    'channel' => 1
]);
$latest = Notification::where('channel', 1)
  ->latest()->take(20)->get();
Notification::query()->delete();
\end{lstlisting}


\subsubsection{Collecting Data}

The \q{Collection}-class wraps arrays with a fluent interface, and it's returned by Eloquent
queries:

\begin{lstlisting}
collect([1, 2, 3, 4])
  ->filter(fn($v) => $v % 2 === 0)  // [2, 4]
  ->map(fn($v) => $v * 10)          // [20, 40]
  ->sum();                          // 60

$heroes->groupBy('clan')->map->count();
$heroes->sortByDesc('level')->pluck('name')->first();
\end{lstlisting}

\noindent Similarly, \q{Str} and \q{Arr} helpers (\q{Str::of(\$text)->trim()->slug()},
\q{Arr::get(\$data, 'user.name')}) simplify operations with strings and nested arrays.


\subsubsection{Communicating with Services}

The \q{Http}-client (a wrapper around Guzzle) sends requests with retries, timeouts, and
parallel execution:

\begin{lstlisting}
$response = Http::withHeaders(
    ['Accept-Language' => 'en-US'])
  ->timeout(10)
  ->retry(3, 200)
  ->asForm()
  ->post('https://example.com/login.php',
    ['login' => $name]);
$response->successful(); // status 2xx
$response->body();       // raw content
$response->json('data'); // decoded value
// Parallel requests
[$chat, $forum] = Http::pool(fn(Pool $pool) => [
  $pool->get('https://example.com/chat'),
  $pool->get('https://example.com/forum'),
]);
\end{lstlisting}


\subsubsection{Processing in Background}

Jobs move long or repeating operations out of the request lifecycle. They are dispatched into a
queue (database, Redis, SQS, or \q{sync}) and processed by a worker (\q{php artisan queue:work}):

\begin{lstlisting}
class ListenStream implements ShouldQueue
{
  use Dispatchable, InteractsWithQueue, Queueable, SerializesModels;

  public $timeout = 60;

  public function __construct(
    public string $generation = '',
    public string $sessionId = '',
  ) {}

  public function handle(): void
  {
    // ... fetch and process new chat messages
    static::dispatch($this->generation, $this->sessionId)
      ->delay(now()->addSeconds(5));
  }
}

ListenStream::dispatch($generation, $sessionId);
\end{lstlisting}

\noindent Events and listeners decouple parts of the application
(\q{event(new MessageReceived(\$msg))}), while the scheduler
(\q{Schedule::job(...)->everyMinute()} in \q{routes/console.php}) replaces a set of cron entries.


\subsubsection{Commanding via Artisan}

Custom commands automate maintenance and debugging tasks. Arguments and options are declared in
the signature: \q{\{name\}} (required argument), \q{\{name?\}} (optional), \q{\{name=default\}}
(with a default value), \q{\{--flag\}} (boolean option), and \q{\{--key=\}} (option with a
value); and a description of each can be added after a colon. Values are read by
\q{\$this->argument()} and \q{\$this->option()}:

\begin{lstlisting}
class StreamListen extends Command
{
  protected $signature = 'app:stream-listen
    {session? : Chat session id (taken from the chat page if omitted)}
    {--url= : Base URL of the game server}
    {--stop : Stop the running listener instead of starting it}';

  protected $description = 'Listen to the stream and process incoming data';

  public function handle(): int
  {
    if ($this->option('stop')) {
      ListenStream::stop();
      $this->info('Listener stopped');
      return self::SUCCESS;
    }

    $sessionId = $this->argument('session');
    if ($sessionId === null) {
      $url = $this->option('url') ?? Defines::URL;
      $htm = (new AppProxyProvider)
        ->boot(rtrim($url, '/') . '/ch/chmain.php');
      $sessionId = ListenStream::sessionIdFrom($htm);
    }

    if ($sessionId === '') {
      $this->error('Session id is not found');
      return self::FAILURE;
    }

    ListenStream::start($sessionId);
    $this->info("Listening to session {$sessionId}");
    return self::SUCCESS;
  }
}
\end{lstlisting}

\noindent The signature also provides a generated help page (\q{--help}), and missing required
arguments are reported before \q{handle()} is called:

\begin{lstlisting}[language=terminal]
# take session from the page
$ php artisan app:stream-listen
# explicit session id
$ php artisan app:stream-listen a1b2c3
$ php artisan app:stream-listen --url=https://x.com
$ php artisan app:stream-listen --stop
$ php artisan app:stream-listen --help
\end{lstlisting}


\subsubsection{Securing Application}

Laravel covers common security concerns by default: CSRF tokens for forms (\q{@csrf}), escaped
output in Blade, parameter binding in queries (protection from SQL injections), password hashing
(\q{Hash::make()}), and encryption by the application key (\q{APP\_KEY}). Authorization is
controlled by gates and policies (\q{Gate::allows('update', \$post)}), and rate limiting
by \q{RateLimiter}-facade or \q{throttle}-middleware.


\subsubsection{Asserting in Tests}

Laravel extends PHPUnit (and Pest) with HTTP, database, and fake assertions. Feature tests boot
the whole application and can send requests to it without a real server
(\q{\$this->get('/age')->assertOk()}):

\begin{lstlisting}
class NotificationChannelTest extends TestCase
{
  // reset the database state for each test
  use RefreshDatabase;

  public function test_social_restrictions_are_preserved(): void
  {
    // Facades can be mocked
    Dialog::shouldReceive('toast')->twice();

    ListenStream::processScript('a("blocked","player",15,15,2,1,54);', false);
    ListenStream::processScript('a("legacy, message","player");', true);
    ListenStream::processScript("c('system message');", false);

    $this->assertDatabaseCount('notifications', 2);
    $this->assertDatabaseHas('notifications', [
      'message' => '<b>player:</b> legacy, message',
      'channel' => 0,
    ]);
  }
}
\end{lstlisting}

\noindent Fakes replace infrastructure in tests (\q{Http::fake()}, \q{Queue::fake()},
\q{Storage::fake()}), and model factories (\q{User::factory()->count(3)->create()}) generate
test data. The \q{php artisan test --parallel} option speeds up the execution (requires
\q{brianium/paratest}-package).
